\begin{Exercise}[label={ex:ripetizioni}]
    Creare un programma che legga da input una sequenza di interi, di lunghezza non nota a priori (al più 10), e terminata da 0. A tal scopo, si
    realizzi una funzione che riceva come parametri di ingresso un vettore e la sua dimensione fisica, e restituisca la dimensione logica del
    vettore. Tale funzione si deve fare carico della fase di lettura e riempimento dell’array. La sequenza può contenere numeri ripetuti.

    Si realizzi una funzione che, ricevuti in ingresso il primo vettore con la sua dimensione logica, ed un secondo vettore con la sua dimensione
    fisica, memorizzi nel secondo vettore tutti gli elementi del primo, ma senza ripetizioni. La funzione restituisca la dimensione logica del
    secondo vettore.

    Si realizzi un main che invochi le funzioni, e che stampi a video l’elenco degli elementi non ripetuti
\end{Exercise}

\begin{Answer}[ref={ex:ripetizioni}]
  	\lstinputlisting[style=cstylesmall]{src/ripetizioni.c}
\end{Answer}